% ============================================================
%  13-conclusion.tex  —  Conclusion générale et perspectives
% ============================================================
\chapterstar{Conclusion Générale et Perspectives}

\begin{chapterintro}
  Ce chapitre clôture le rapport en synthétisant les réalisations du projet
  SolarEase, en dressant un bilan critique et en exposant les perspectives
  d'évolution à court et moyen terme.
\end{chapterintro}

\section*{Synthèse des réalisations}

Le projet \textbf{SolarEase} a permis de concevoir et de réaliser une
plateforme web complète répondant aux besoins des sociétés d'installation
de panneaux solaires. Organisé en sprints Kanban successifs, le travail
accompli couvre~:

\begin{itemize}
  \item \textbf{Chapitre~3 (Sprints 1 et 2)~:} socle technique et gestion
    métier~: authentification sécurisée (inscription client OTP, comptes
    installateurs par l'admin, JWT, Gateway), gestion des clients et
    projets (CRUD, cycle de vie 6 statuts, tableau de bord KPI)~;
  \item \textbf{Chapitre~4 (Sprints 3, 4 et 5)~:} cœur métier et
    collaboration~: moteur de dimensionnement photovoltaïque (données PVGIS
    réelles, modèle Night Panel, analyse financière sur 25 ans), approche
    RAG \textbf{hybride} pour les recommandations IA (sélection
    déterministe sur le catalogue société + base vectorielle
    \texttt{pgvector} + Ollama), génération automatique de rapports PDF
    professionnels, suivi terrain avec mises à jour de progression,
    messagerie temps réel (WebSocket), centre de notifications, espace
    client dédié avec stepper de suivi, gestion des demandes solaires,
    devis avec acceptation client et facturation automatique~;
  \item \textbf{DevOps (Chapitre~5)~:} industrialisation complète de la
    chaîne logicielle backend~: containerisation multi-stage des 4
    microservices, pipeline CI/CD GitHub Actions (build, tests, scans
    Trivy/Gitleaks, publication d'images sur GHCR), stack d'observabilité
    (Spring Actuator + Prometheus + Grafana + Loki), qualité (JaCoCo),
    Infrastructure as Code (Kubernetes + Kustomize, overlays dev/prod),
    déploiement automatique Railway sur les branches \texttt{develop} et
    les tags \texttt{v*.*.*}. Le frontend React/Vite reste un projet
    léger livré en continu sur Vercel.
\end{itemize}

\section*{Bilan des objectifs atteints}

\begin{table}[H]
  \caption{Bilan des objectifs du projet SolarEase}
  \label{tab:bilan}
  \centering
  \begin{tabularx}{\textwidth}{|L{5cm}|C{2cm}|X|}
    \hline
    \rowcolor{figblue}
    \textcolor{white}{\textbf{Objectif}} &
    \textcolor{white}{\textbf{Atteint}} &
    \textcolor{white}{\textbf{Commentaire}} \\
    \hline
    Authentification sécurisée JWT+OTP & \textcolor{green!70!black}{\checkmark} Oui & Entièrement fonctionnel \\
    \hline
    Gestion clients et projets & \textcolor{green!70!black}{\checkmark} Oui & CRUD + cycle de vie \\
    \hline
    Moteur de dimensionnement & \textcolor{green!70!black}{\checkmark} Oui & PVGIS + fallback \\
    \hline
    Innovation Night Panel & \textcolor{green!70!black}{\checkmark} Oui & Formules validées \\
    \hline
    Analyse financière 25 ans & \textcolor{green!70!black}{\checkmark} Oui & ROI, payback, cash-flow \\
    \hline
    Génération rapport PDF & \textcolor{green!70!black}{\checkmark} Oui & Téléchargement direct \\
    \hline
    Recommandation IA & \textcolor{green!70!black}{\checkmark} Oui & Ollama en fallback contrôlé \\
    \hline
    Architecture microservices & \textcolor{green!70!black}{\checkmark} Oui & 4 services + Docker Compose \\
    \hline
    Pipeline CI/CD automatisé & \textcolor{green!70!black}{\checkmark} Oui & GitHub Actions + GHCR \\
    \hline
    Observabilité (métriques + logs) & \textcolor{green!70!black}{\checkmark} Oui & Prometheus + Grafana + Loki \\
    \hline
    Infrastructure as Code & \textcolor{green!70!black}{\checkmark} Oui & Kustomize (base + overlays) \\
    \hline
    Déploiement multi-environnement & \textcolor{green!70!black}{\checkmark} Oui & Railway staging + prod \\
    \hline
  \end{tabularx}
\end{table}

\section*{Bilan critique}

\textbf{Points forts~:}
\begin{itemize}
  \item Architecture moderne, découplée et évolutive~;
  \item Données de production basées sur PVGIS (source officielle européenne)~;
  \item Innovation significative avec le mode Night~Panel et la couverture nocturne~;
  \item Expérience utilisateur professionnelle avec deux espaces distincts~;
  \item Résilience assurée par les mécanismes de fallback (PVGIS et IA).
\end{itemize}

\textbf{Limites identifiées~:}
\begin{itemize}
  \item Dépendance partielle à des services externes (PVGIS, Ollama)~;
  \item Paramètres de calibration (prix kWh, taux dégradation) à valider avec
    chaque société partenaire~;
  \item Absence de tests de charge exhaustifs en conditions de production.
\end{itemize}

\section*{Perspectives d'évolution}

À court terme~:
\begin{itemize}
  \item Intégration d'un \textbf{tableau de bord analytique avancé} avec
    visualisations des économies réalisées~;
  \item Enrichissement de la \textbf{base de connaissances IA} avec des
    données spécifiques au marché tunisien~;
  \item Ajout d'un \textbf{module de signature électronique} pour les devis.
\end{itemize}

À moyen terme~:
\begin{itemize}
  \item \textbf{Architecture multi-tenant} pour servir plusieurs sociétés
    d'installation depuis une même instance~;
  \item \textbf{Application mobile} React Native pour les commerciaux terrain~;
  \item Mise en place d'\textbf{Alertmanager} couplé à Slack/Teams pour les
    alertes critiques sur les métriques Prometheus déjà collectées~;
  \item Tests \textbf{end-to-end} automatisés (Playwright) intégrés au
    pipeline GitHub Actions~;
  \item Adoption d'un \textbf{service mesh} (Istio ou Linkerd) pour la
    résilience inter-services en production~;
  \item Connexion aux systèmes de \textbf{monitoring en temps réel} des
    installations existantes.
\end{itemize}

\vspace{0.8cm}
\begin{center}
  \begin{tcolorbox}[colback=isetblue!8, colframe=isetblue,
      width=0.9\textwidth, halign=center, fontupper=\itshape]
    SolarEase démontre qu'il est possible de développer, en solo, une
    plateforme professionnelle complète en appliquant rigoureusement les
    bonnes pratiques d'ingénierie logicielle~: architecture microservices,
    sécurité centralisée, tests automatisés et déploiement conteneurisé.
    Ce projet constitue une base solide et extensible pour une mise en
    production réelle dans le secteur des énergies renouvelables.
  \end{tcolorbox}
\end{center}

\cleardoublepage
